\section{Numerical Methods}
Numerical methods are ways to approximate the solution to equations using a repeatable procedure, usually designed for computers to use efficiently.

In an exam, you should assume that you are to solve a question exactly unless specifically told otherwise, e.g. ``by using the trapezium rule, estimate the integral.''

In fact, it is always preferred to get an exact solution by solving mathematically rather than estimating using one of these techniques, and these should be reserved for when that is not possible or feasible.

\subsection{The Trapezium Rule}\label{sec:traprule}

The goal of the trapezium rule is to estimate the area under the curve\footnote{Actually we are estimating the integral, which is a signed area since the integral of a negative function is also negative. This is left here to assist intuition.} as in the shaded area of Figure \ref{fig:shaded-function}, i.e. a way to estimate an integral, using only knowledge of the function we are trying to integrate at a few points.

\begin{figure}[H]
\centering
\incfig[0.75]{shaded function}
\caption{A function shaded below the curve}
\label{fig:shaded-function}
\end{figure}

One approach is to:

\begin{enumerate}
    \item Choose a bunch of $x$ points we want to know the value of the function at.
    \item Find what $f(x)$ is at each point chosen. 
    \item Connect the $f(x)$ points with straight lines, and draw lines down to the $x$ axis from each $f(x)$ point found. 
    This will make a sequence of trapeziums, where the ``top'' of the trapezium is approximating the function curve with straight lines.
    This can be seen in Figure \ref{fig:shaded-function-with-trapz}.
    \item Calculate the area of each trapezium and add them together.
\end{enumerate}

\begin{figure}[H]
\centering
\incfig[0.75]{shaded function with trapz}
\caption{A function shaded below the curve, with trapeziums drawn, equal width.}
\label{fig:shaded-function-with-trapz}
\end{figure}

Recall that for a trapezium with side lengths $a$ and $b$, and height $h$, the area of a trapezium is the following.

\[ \text{Area of trapezium} = \frac{h(a+b)}{2}\]

This can be proved by formalising the geometric construction seen in Figure \ref{fig:trapezium-geometric}, creating a copy of the trapezium, rotating it $180^\circ$, and then gluing it on the side of the original trapezium.
It can be proved\footnote{You would approach this by labelling the angles in the outer triangles, and using the fact that they must add up to $180^\circ$ to find that the construction does indeed create a parallelogram.} 
that this creates a parallelogram with width $a+b$ and height $h$, which has a total area of $h(a+b)$. Then since we duplicated our trapezium, this must be twice the area of the original trapezium.


\begin{figure}[H]
\centering
\incfig[0.75]{trapezium}
\caption{A trapezium glued to itself to form a parallelogram.}
\label{fig:trapezium-geometric}
\end{figure}


So what are the areas of the trapeziums in our problem? For convenience, we will space the $x$ points out evenly, with a distance of $h$ between each one. 
As in, $x_1 = x_0 + h$, $x_2 = x_1 + h$, and so on for how many $x$ points we chose.

Firstly, let's label the diagram with the widths and heights of the trapeziums, as in Figure \ref{fig:shaded-function-with-trapz-labelled}.

\begin{figure}[H]
\centering
\incfig[0.75]{shaded function with trapz labelled}
\caption{A function shaded below the curve, with trapeziums drawn, equal width.}
\label{fig:shaded-function-with-trapz-labelled}
\end{figure}

From this figure, we can notice that each trapezium has a side length of $f(x_i)$ and $f(x_{i+1})$, and a height of $h$.

Therefore, the area of one of the trapeziums must be of the form:

\[ \text{Area of trapezium }i = \frac{h(f(x_i) + f(x_{i+1}))}{2}\]

To approximate the total area of $f(x)$ between two points, we add them all together to get (starting from $x_0$ and ending at $x_n$):

\begin{align*}
    \int_{x_0}^{x_n}f(x) dx \approx& \frac{h\big(f(x_0) + \textcolor{red}{f(x_{1})}\big)}{2} \tag{first trapezium} \\
     &+ \frac{h\big(\textcolor{red}{f(x_{1})} + \textcolor{teal}{f(x_{2})}\big)}{2} \tag{second trapezium}\\
     &+ \frac{h\big(\textcolor{teal}{f(x_{2})} + f(x_{3})\big)}{2} \tag{third trapezium}\\
     &+ \dots \\
     &+ \frac{h\big(f(x_{n-2}) + \textcolor{olive}{f(x_{n-1})}\big)}{2} \tag{second-to-last trapezium}\\
     &+ \frac{h\big(\textcolor{olive}{f(x_{n-1})} + f(x_{n})\big)}{2} \tag{last trapezium} 
\end{align*}

Similar terms have been colour-coded. In fact, the right-hand-side can be simplified drastically by collecting the similar terms. 
You can see that the term at the end of the line appears at the start of the next line, which means that there are repeat terms for every line except 
the first or the last one. Let's expand and collect the terms together to see what we get:
\begin{align*}
    \int_{x_0}^{x_n}f(x) dx \approx \frac{h}{2}\bigg(& f(x_0) + \textcolor{red}{f(x_{1})} \tag{pulling $h/2$ out of the full equation} \\
     &+ \textcolor{red}{f(x_{1})} + \textcolor{teal}{f(x_{2})} \\
     &+ \textcolor{teal}{f(x_{2})} + f(x_{3}) \\
     &+ \dots \\
     &+ f(x_{n-2}) + \textcolor{olive}{f(x_{n-1})} \\
     &+ \textcolor{olive}{f(x_{n-1})} + f(x_{n}) \bigg) \\
     = \frac{h}{2} \bigg( & f(x_0) + 2\textcolor{red}{f(x_{1})}+ 2\textcolor{teal}{f(x_{2})}  + \dots + 2\textcolor{olive}{f(x_{n-1})}+ f(x_n)\bigg) \tag{collecting terms}\\
     = \frac{h}{2} \bigg( & f(x_0) + f(x_n) + 2\big( \textcolor{red}{f(x_{1})}+ \textcolor{teal}{f(x_{2})} + \dots + \textcolor{olive}{f(x_{n-1})}\big) \bigg)  \tag{rearranging}
\end{align*}

This gives us the trapezium rule from the formula book:
\[ \int_{x_0}^{x_n}f(x) dx  \approx  \frac{h}{2} \bigg(  f(x_0) + f(x_n) + 2\big( f(x_{1})+ f(x_{2}) + \dots + f(x_{n-1})\big) \bigg) \]

Now the question is: what is $h$ in terms of $n$?

If we know what the start and end points are that we are integrating between, and we know how many points we are going to use, we can figure out what $h$ is.

To split the line-segment between 
$x_0$ and $x_n$ into $n$ segments, we need to set $h$ to be 

\[ h = \frac{x_n - x_0}{n}\]

So that 
\[ x_{i+1} = x_i + h \]

Note that since we are counting from $x_0$ to $x_n$, we actually have $(n+1)$ points, but $n$ segments between them.

\begin{example}
    By splitting the interval between $0$ and $3$ with $3$ equally spaced segments, use the trapezium rule to estimate the integral:

    \[ \int_0^3 x^2 dx\]

    Then compare this with the answer you get from integrating it directly.
\end{example}

\begin{solution}
    The $3$ equally spaced segments will have $h = \frac{3 - 0}{3} = \frac{3}{3} = 1$. That means we have:
    \begin{align*}
        x_0 &= 0 \\
        x_1 &= x_0 + h = 1 \\
        x_2 &= x_1 + h = 2 \\
        x_3 &= x_2 + h = 3
    \end{align*}

    Setting $f(x) = x^2$, let's compute each value as a table:

    \begin{center}
    \begin{tabular}{c|cccc}
    $x$        & 0 & 1 & 2 & 3 \\
    \hline
    $f(x)=x^2$ & 0 & 1 & 4 & 9
    \end{tabular}
    \end{center}

    Now let's plug this into the trapezium rule:

    \begin{align*}
    \int_{x_0}^{x_n}f(x) dx  &\approx  \frac{h}{2} \bigg(  f(x_0) + f(x_n) + 2\big( f(x_{1})+ f(x_{2}) + \dots + f(x_{n-1})\big) \bigg) \\
    \therefore \int_{0}^{3}x^2 dx  &\approx  \frac{1}{2} \bigg(  0 + 9 + 2\big( 1 + 4 \big) \bigg) \\
    &= \frac{1}{2} \bigg(  9+ 2\big( 5 \big) \bigg) \\
    &= \frac{19}{2}\\
    &= 9.5
    \end{align*} 

    Comparing this to computing the integral directly:

    \begin{align*}
    \int_{x_0}^{x_n}f(x) dx  &=  \int_0^3 x^2 dx \\
     &= \left[\frac{1}{3}x^3\right]_{x=0}^{x=3}\\ 
     &= \frac{1}{3}3^3 \\
     &= 3^2 \\
     &= 9
    \end{align*} 
    Comparing it with the trapezium rule, we were only out by $0.5$ -- the trapezium rule with $4$ equally spaced segments got $9.5$ as the integral whereas the actual answer was $9$. Pretty close!
    \qed
\end{solution}

In the previous question, we found that the trapezium rule overshot the real answer. If you draw out the function with the trapeziums, you can see this visually -- the trapeziums include sections above the curve since it curves upwards. 
In fact, you can usually anticipate whether the trapezium rule will over or undershoot the actual integral 
by looking at the section which does not overlap between the trapeziums and the function itself, look at Figure \ref{fig:shaded-function-with-trapz} again to see which parts the trapeziums include more area than the function, and which parts the trapeziums include less area than the function. Notice the patterns in the overlap when the function is curving in the same direction.

\begin{example}
    You are given the following grid for of values for a function.

        \begin{center}
    \begin{tabular}{c|cccc}
    $x$        & -1 & 0.25 & 1.5 & 2.75 \\
    \hline
    $f(x)$ & -1 & 4 & 6 & 5
    \end{tabular}
    \end{center}

    By using the Trapezium rule, estimate the integral of the curve between $x=-1$ and $x=2.75$.

    If you are told that the curve's derivative is always decreasing, what can you say about whether the estimate for the area will be under or over the true area of $f(x)$?
\end{example}
\begin{solution}
    First we need to determine $h$. Since $h$ is the distance between two $x$ points, we can deduce that it should be:
    \[ h= 0.25 - (-1) = 1.25\]

    And we can verify that each point is $1.25$ apart, meaning that $h=1.25$ is the distance between each point, and we can use the trapezium rule.
    Now let's plug this into the trapezium rule:

    \begin{align*}
    \int_{x_0}^{x_n}f(x) dx  &\approx  \frac{h}{2} \bigg(  f(x_0) + f(x_n) + 2\big( f(x_{1})+ f(x_{2}) + \dots + f(x_{n-1})\big) \bigg) \\
    \therefore \int_{-1}^{2.75}f(x) dx  &\approx  \frac{1.25}{2} \bigg(  -1 + 5 + 2\big( 4 + 6 \big) \bigg) \\
    &= \frac{1.25}{2} \bigg(  4+ 2\big( 10 \big) \bigg) \\
    &= \frac{1.25 \times 24}{2}\\
    &= 15
    \end{align*} 

    Therefore the trapezium rule gives us an estimate of the area as $15$.

    For the second part of the question, if we draw a sketch of the curve, we know that the slope must be going from steep to less steep since the question told us the derivative is decreasing.

    \begin{figure}[H]
    \centering
    \incfig[0.75]{shaded function 2}
    \caption{A function shaded below the curve, with trapeziums drawn.}
    \label{fig:shaded-function-2}
    \end{figure}

    Now it becomes clear that the trapeziums are always under the curve. When the curve is above zero, this means that the trapezium area is less than the curve's area. What about when the function 
    is below zero? In this case, the trapezium area is covering more space than the function is. But since it is negative, this should contribute more ``negative area''. However since we have not marked zero itself in the table, the trapezium calculation we are computing has negative side (!), which affects the area in this case. 
    
    The actual ``trapezium'' we are computing for the first trapezium is:

    \[ \frac{h(-1 + 5)}{2}\]

    This makes it unclear whether the area in this case is less than or more than the actual function, at least for the geometrical approach we have been using, but it turns out the trapezium rule still underestimates the true integral in this case too, as long as the curve's derivative is decreasing\footnote{If we instead think of the straight line segment itself between $x_0$ and $x_1$ as a function, say $L(x)$, it follows due to the decreasing derivative that $f(x)\geq L(x)$ on this interval, and the ``trapezium area'' we have been using is given as $\int_{x_0}^{x_1}L(x)dx$. It then follows that $\int_{x_0}^{x_1}f(x)dx \geq\int_{x_0}^{x_1}L(x)dx$, which shows that the trapezium rule still undershoots the true value in this case.}.
    \qed
\end{solution}


\subsection{The Newton-Raphson Method}
The Newton-Raphson method has the goal of improving an estimate for a root of a function, i.e. finding a point $x$ where some function $f(x) = 0$, given that we know a point close to the root already.

The curve pictured in Figure \ref{fig:function-with-zeros-labelled} has three roots. We could try to solve $f(x) = 0$ by rearranging for $x$, but this is not always possible to do.

Instead, we can try to get \emph{close} to a root. The easiest way to do so is to find two points $a$ and $b$ where the function changes sign between the two points. If the function goes from positive to negative, or from negative to positive, 
we know that it must be zero at some point between. This will give us a point close to the root.

\begin{figure}[H]
\centering
\incfig[0.75]{function with zeros labelled}
\caption{A function with zeros labelled.}
\label{fig:function-with-zeros-labelled}
\end{figure}

Once we have a point close to the root, how can we improve our guess? One method is to notice that near the root, if we zoom in far enough, the function looks like a straight line. Because of this, 
we might notice that if we were to draw the tangent to the curve at our point, following the tangent until we hit the $x$ axis should land us closer to the root.

This is the Newton-Raphson method: 

\begin{enumerate}
    \item Find a point close to the root 
    \item Find the tangent to the curve at that point 
    \item Find the point on the $x$ axis in which the tangent meets the $x$ axis 
    \item Repeat steps 2-3 until you are satisfied with getting near the root
\end{enumerate}

This is seen pictured in Figure \ref{fig:function-with-newton-iterations-labelled}, where you can see the tangents leading us closer to the red point we desire to land at.

\begin{figure}[H]
\centering
\incfig[0.75]{function with zeros labelled_newton}
\caption{A function with Newton iterations labelled around one root.}
\label{fig:function-with-newton-iterations-labelled}
\end{figure}

So how can we turn the above instructions into maths?

Firstly, we want to find the gradient of our function $f$ at an initial point $x_0$. This is given by $f'(x_0)$, as in by differentiating $f$ and plugging in $x_0$.

Then given this gradient, we want to find the equation of the straight line which has this gradient, and passes through the point $f(x_0)$, which is exactly the properties we want from 
the tangent to a curve! It needs to ``touch'' the curve and have the slope the curve has at that point.

We can write a general straight line as $y = mx + c$, where $m$ is the gradient, and $c$ is the $y$ intersect. But we know that at $x_0$, our tangent line passes through $f(x_0)$ with gradient $m=f'(x_0)$.

We can write this tangent observation as an equation:

\begin{align*}
    y &= f'(x) x +c \tag{our general equation for a straight line}\\
    \therefore f(x_0) &= f'(x_0)x_0 + c \tag{using our observations above}\\
    \therefore c &= f(x_0) - f'(x_0)x_0 \tag{rearranging}\\
    \therefore y &= f'(x_0)x + f(x_0) - f'(x_0)x_0 \tag{subbing into original equation}
\end{align*}

Now that we have found the equation of a general tangent line to $f(x)$ at the point $x_0$, we can find where this crosses the $x$-axis.
To find this, we simply need to set $y=0$ and solve for $x$:

\begin{align*}
    0 &= f'(x_0)x + f(x_0) - f'(x_0)x_0 \tag{setting $y=0$}\\
    \therefore 0 &= x + \frac{f(x_0)}{f'(x_0)} - x_0 \tag{dividing by $f'(x_0)$}\\
    \therefore x &= x_0 - \frac{f(x_0)}{f'(x_0)} 
\end{align*}

We can convert this into an interative process by considering the point on the left side to be the ``next'' point, and then run through the same process again but swapping $x_0$ with the next point. This gives the 
\textbf{Newton-Raphson Iteration for solving $f(x)=0$:}

\[ x_{n+1} = x_n - \frac{f(x_n)}{f'(x_n)}\]

Notice that we divide by $f'(x_n)$, and therefore this formula does not work if the derivative becomes zero at some point.

Also notice that, since this method assumes that our function ``looks like'' a straight line when you zoom in, and that we chose an initial point in this general region of straight-ness. If our initial guess of $x_0$ is not close enough to the actual root, then this method may not approach the root even after many iterations.

\begin{example}
    Consider $f(x) = x^5 + x^3 - 1$. Show that there is a root between $0$ and $1$.
    By choosing $x_0 = 1$, compute two Newton-Raphson iterations to get an approximation of a root to this function.
\end{example}
\begin{solution}
    Firstly, to check that there is a root between $0$ and $1$, we can plug in $0$ and $1$ to see if the function crosses zero in-between the two numbers. If one is negative and the other is positive, then we know this must have happened.

    First plug in $0$:
    \[ f(0) = 0^5+0^3-1 = -1\]

    Now plug in $1$:
    \[ f(1) = 1^5 + 1^3 - 1 = 1\]

    Since at $0$ the function is negative and at $1$ the function is positive, we know it crossed zero somewhere in-between.

    Now for Newton-Raphson, we need to find the derivative of the function.

    \[f'(x) = 5x^4 + 3x^2 \]

    Now plug into the formula starting with $x_0 = 1$:
    \[ x_{n+1} = x_n - \frac{f(x_n)}{f'(x_n)}\]
    \begin{align*}
        x_1 &= x_0 - \frac{f(x_0)}{f'(x_0)} \\
         &= 1 - \frac{f(1)}{f'(1)}\\
         &= 1 - \frac{1}{5 + 3 } \\
         &= 1 - \frac{1}{8}\\
         &= \frac{7}{8}
    \end{align*}

    For the second iteration, we use that $x_1 = 6/7$ to find:
    \begin{align*}
        x_2 &= x_1 - \frac{f(x_1)}{f'(x_1)} \\
            &= \frac{7}{8} - 
            \frac{\left(\frac{7}{8}\right)^5 + \left(\frac{7}{8}\right)^3 - 1}
            {5\left(\frac{7}{8}\right)^4 + 3\left(\frac{7}{8}\right)^2} \\
            &\approx 0.84003
    \end{align*}

    Therefore, after two Newton-Raphson iterations, the root is approximately
    \[
        x \approx 0.84003.
    \]
    \qed
\end{solution}
